\clearpage
\thispagestyle{empty}
\begin{center}
    \bfseries\fontsize{14}{16}\selectfont ABSTRACT
\end{center}
\vspace{16pt}

\noindent
The \textit{Capacitated Vehicle Routing Problem} (CVRP) is a classical combinatorial
optimization problem in which a minimum-cost set of routes must be designed to serve
customers with known demands by means of a homogeneous fleet of capacity-constrained
vehicles. In Report~No.~1 the problem was addressed through exact mathematical
programming methods, whose computational effort grows sharply with instance size. This
report tackles the CVRP from the standpoint of \textit{metaheuristics}, implementing
\textit{Particle Swarm Optimization} (PSO) as an approximate method able to obtain
high-quality solutions without a certificate of optimality, and contrasts both
approaches over a common set of instances. To adapt PSO, a continuous-domain
algorithm, to the combinatorial search space of the CVRP, the SR-1 representation
\parencite{ai2009particle} was implemented, which decodes each particle into a customer
priority list and vehicle reference points, building feasible routes through
least-cost insertion and 2-opt local improvement. PSO hyperparameters were calibrated
with Optuna (60 trials) on tuning instances distinct from the evaluation set. The
method was evaluated with 31 independent repetitions over nine CVRPLIB instances: the
five from Report~1 (for direct comparison) and four additional, larger instances beyond
the reach of the exact method. On the five shared instances, PSO achieved an average GAP
between $0.03\%$ and $10.07\%$ with respect to the optimum, in sub-second runtimes,
versus an exact method that took up to $4.7$ hours without certifying optimality on the
largest shared instance. On the four additional instances, the average GAP ranged from
$15.25\%$ to $26.84\%$, with the maximum on the largest instance. The comparison reveals
a complementarity between
both approaches: the exact method certifies optimality while it remains tractable, but
loses tractability abruptly beyond a certain size; PSO does not certify optimality, but
keeps bounded runtimes and solves instances considerably larger than those reachable by
the exact method.

\vspace{10pt}
\noindent\textbf{Keywords:} CVRP; combinatorial optimization; metaheuristics;
\textit{Particle Swarm Optimization}; PSO; CVRPLIB.
